\documentclass[11pt]{article}
\usepackage[utf8]{inputenc}
\usepackage{amsmath,amssymb}
\usepackage{hyperref}
\title{Efficient Serverless Computing Scheduling: A Resource-Aware Approach}
\author{Anonymous}
\date{}
\begin{document}
\maketitle

\begin{abstract}
This paper studies Serverless Computing Scheduling. Grounded in a real, citable literature \cite{ref1,ref2,ref3,ref4,ref5,ref6,ref7,ref8},
we propose a focused method and report \textbf{illustrative} experiments.
All references below are real and verifiable (OpenAlex/arXiv); the reported
numbers are simulated to illustrate intended behaviour.
\end{abstract}

\section{Introduction}
Serverless Computing Scheduling is an active research area. This work targets a concrete gap identified from
the recent literature and contributes a method together with an evaluation protocol.

\section{Related Work (real literature)}
The following works, retrieved from public bibliographic APIs, frame our study:
\begin{itemize}
\item Walia et al. (2023) — "AI-Empowered Fog/Edge Resource Management for IoT Applications: A Comprehensive Review, Research Challenges, and Future Perspectives", IEEE Communications Surveys \& Tutorials (cited 249).
\item Li et al. (2022) — "Serverless Computing: State-of-the-Art, Challenges and Opportunities", IEEE Transactions on Services Computing (cited 190).
\item Rausch et al. (2020) — "Optimized container scheduling for data-intensive serverless edge computing", Future Generation Computer Systems (cited 180).
\item Xie et al. (2021) — "When Serverless Computing Meets Edge Computing: Architecture, Challenges, and Open Issues", IEEE Wireless Communications (cited 155).
\item Adhikari et al. (2019) — "A Survey on Scheduling Strategies for Workflows in Cloud Environment and Emerging Trends", ACM Computing Surveys (cited 149).
\item Wang et al. (2023) — "Deep Reinforcement Learning-based scheduling for optimizing system load and response time in edge and fog computing environments", Future Generation Computer Systems (cited 144).
\end{itemize}

\section{Method}
We reduce the dominant computational cost via adaptive resource allocation, activating only the components needed per input. We instantiate this idea for Serverless Computing Scheduling and analyse its cost/quality trade-off.

\section{Experiments (Illustrative)}
\begin{center}
\begin{tabular}{lcc}
System & Primary Metric & Efficiency \\ \hline
Strong baseline & 0.812 & 1.00$\times$ \\
Ablation & 0.828 & 0.92$\times$ \\
\textbf{Ours} & \textbf{0.851} & \textbf{0.71$\times$}
\end{tabular}
\end{center}
\emph{Metrics simulated to illustrate the intended effect; not measured.}

\section{Conclusion}
We connected a real literature to a concrete method for Serverless Computing Scheduling. Future work will
validate the illustrative results with real experiments.

\begin{thebibliography}{9}
\bibitem{ref1} Guneet Kaur Walia, Mohit Kumar, Sukhpal Singh Gill. AI-Empowered Fog/Edge Resource Management for IoT Applications: A Comprehensive Review, Research Challenges, and Future Perspectives. \emph{IEEE Communications Surveys \& Tutorials}, 2023. DOI:10.1109/comst.2023.3338015
\bibitem{ref2} Yongkang Li, Yanying Lin, Yang Wang et al.. Serverless Computing: State-of-the-Art, Challenges and Opportunities. \emph{IEEE Transactions on Services Computing}, 2022. DOI:10.1109/tsc.2022.3166553
\bibitem{ref3} Thomas Rausch, Alexander Rashed, Schahram Dustdar. Optimized container scheduling for data-intensive serverless edge computing. \emph{Future Generation Computer Systems}, 2020. DOI:10.1016/j.future.2020.07.017
\bibitem{ref4} Renchao Xie, Qinqin Tang, Shi Qiao et al.. When Serverless Computing Meets Edge Computing: Architecture, Challenges, and Open Issues. \emph{IEEE Wireless Communications}, 2021. DOI:10.1109/mwc.001.2000466
\bibitem{ref5} Mainak Adhikari, Tarachand Amgoth, Satish Narayana Srirama. A Survey on Scheduling Strategies for Workflows in Cloud Environment and Emerging Trends. \emph{ACM Computing Surveys}, 2019. DOI:10.1145/3325097
\bibitem{ref6} Zhiyu Wang, Mohammad Goudarzi, Mingming Gong et al.. Deep Reinforcement Learning-based scheduling for optimizing system load and response time in edge and fog computing environments. \emph{Future Generation Computer Systems}, 2023. DOI:10.1016/j.future.2023.10.012
\bibitem{ref7} Qinqin Tang, Renchao Xie, F. Richard Yu et al.. Distributed Task Scheduling in Serverless Edge Computing Networks for the Internet of Things: A Learning Approach. \emph{IEEE Internet of Things Journal}, 2022. DOI:10.1109/jiot.2022.3167417
\bibitem{ref8} Kostis Kaffes, Neeraja J. Yadwadkar, Christos Kozyrakis. Centralized Core-granular Scheduling for Serverless Functions. \emph{}, 2019. DOI:10.1145/3357223.3362709
\end{thebibliography}
\end{document}
